\section{Introduction}

The proliferation of autonomous multi-agent systems across critical infrastructure—financial markets, agricultural operations, healthcare logistics, and defense systems—has exposed a structural gap in the architecture of accountability. Modern autonomous agents are designed to decompose complex goals into recursive sub-tasks, delegate authority across hierarchies of machine actors, and execute without ongoing human oversight. Yet the prevailing architecture provides no mandatory mechanism for routing exceptional conditions to a human decision-maker when all machine resolution paths fail. The result is a class of failure modes in which autonomous systems proceed to act without accountability, not because they are malicious, but because the architecture contains no path back to human consciousness when it is genuinely needed.

This problem is not theoretical. When a recursive multi-agent system encounters an unanticipated condition—a sensor reading that violates all trained distributions, a delegation chain that exceeds its scope boundary, a safety envelope that is breached in a novel way—the system's behavioral options are typically limited to local error handling, retry loops, or shutdown. Local error handling fails when the condition exceeds the local actor's representational capacity. Retry loops fail when the condition is not a transient noise artifact but a genuine structural mismatch. Shutdown is the correct response in only a subset of failure cases; many conditions require judgment that the system was not designed to exercise. In the absence of a mandatory escalation path, the system does not stop. It continues, extrapolating from a context that no longer maps to the world it is acting upon. The outcome is actions taken with no accountable human in the loop—actions that may be irreversible, physically hazardous, or legally consequential.

We term the architectural element that addresses this failure mode the \textbf{Bailout Protocol}. The Bailout Protocol is not an error-handling subroutine. It is not a rollback mechanism or a circuit breaker in the conventional sense. It is a structurally mandatory escalation path that is embedded in every node of a recursive autonomous system and that bypasses all machine actors in the resolution chain when the node's local \purpose and \bounds mechanisms cannot resolve a condition with certainty. The protocol routes the exception upward through the immutable parent chain until it reaches a \textbf{Human Accountability Anchor}—a designated human being who is architecturally bound to receive the escalation, to engage with it, and to provide a resolution decision. The system fails upward into human consciousness. It never fails downward into autonomous continuation without accountability.

The core insight underlying the Bailout Protocol is that accountability cannot be a policy constraint—it must be an architectural fallback. Policy constraints are effective under normal operating conditions and under adversarial attack models that the policy was designed to address. They are ineffective under conditions of novelty: situations that were not anticipated in the policy design, edge cases that fall between existing categories, or systemic conditions that emerge from the interaction of multiple subsystems in ways that no single policy contemplated. When a system encounters genuine novelty, the policy layer has no valid response, and without an architectural fallback, the system proceeds on the basis of an invalid context. Accountability architecture must survive the failure of every other accountability mechanism. It must be the element that remains when every other constraint has been exhausted—not because it is the most sophisticated element, but because it is the one that is least likely to be bypassed when conditions become extreme.

The \bxthree framework provides the architectural context for the Bailout Protocol. In \bxthree, every node in a recursive system carries a \purpose layer that defines what it is for, a \bounds engine that defines the limits of permissible action, and a \fact layer that grounds the node's representations in observable reality. Together, these three components define the node's operating context. When a condition arises that the node cannot resolve within its \purpose, that violates its \bounds, or that its \fact layer cannot ground in verifiable state, the node enters an exception condition. In conventional architectures, the exception condition would be handled locally or propagated through a machine-to-machine escalation path. In the \bxthree architecture, the exception condition triggers the Bailout Protocol, which bypasses machine escalation entirely and routes to the Human Accountability Anchor. This behavior is not configurable by the node. It is a structural property of the architecture, embedded in the node's Worksheet, and cannot be disabled, delayed, or rerouted by any machine actor in the system.

The failure of accountability under recursive autonomous conditions has been observed in limited ways in prior work. Arnold and Bommarito's analysis of autonomous system safety guarantees identifies the absence of human oversight in recursive delegation as a critical failure mode, though their proposed mitigations focus on policy-level constraints that are subject to the novelty limitations discussed above.\footnote{See Arnold and Bommarito's framework for autonomous system safety in multi-agent environments.} The OODA loop literature, originating with Boyd, recognizes the necessity of a human decision-maker in the outer loop of autonomous systems operating at speed, but does not address the architecture by which a human is guaranteed to receive exceptions that machine actors cannot resolve.\footnote{Boyd's OODA loop framework establishes the importance of human engagement in the decision cycle but does not specify the escalation architecture for machine-only resolution failure.} More recent work on multi-agent robustness has identified cascading failure modes in autonomous systems but has focused primarily on mitigation strategies within the machine domain, without addressing the specific case where all machine-domain mitigations have been exhausted and a human must be reached.\footnote{See multi-agent robustness literature on cascading failures and systemic risk in autonomous networks.} The Bailout Protocol contributes a specific architectural solution: a mandatory, non-disableable, upward-escalating exception path that is structurally embedded at every node and that is triggered by the exhaustion of machine resolution capacity rather than by any particular policy violation.

This paper proceeds as follows. In Section 2, we formally define the Bailout Protocol and its triggering conditions, distinguishing between three trigger classes: capability boundary exhaustion, safety envelope violation, and accountability boundary violation. In Section 3, we describe the architecture of the human accountability anchor and the immutable parent chain that ensures escalation reaches it. In Section 4, we analyze the operational properties of the protocol, including its behavior under adversarial conditions, its performance under normal operation, and its behavior when the human anchor is unavailable. In Section 5, we present the Worksheet specification that embeds the Bailout Protocol in every node. In Section 6, we demonstrate the protocol's behavior through a case study in agricultural autonomous systems, where the failure of a sensor-driven irrigation node triggers a bailout that is routed to a human operator despite the node's complete disconnection from the central system. In Section 7, we discuss the implications of the Bailout Protocol for regulatory frameworks, including the EU AI Act's requirements for human oversight in high-risk AI systems. In Section 8, we offer conclusions and describe open research questions.